\chapter{Design}
[IMAGE PLACEHOLDER: methodology_flow.png]
The design isolates the core variables.

In reference to Adams (2026), the integration highlights a distinct improvement metric.


Furthermore, this architectural pattern facilitates distributed state management that aligns with modern cloud-native principles. 
By effectively decoupling the data ingestion from processing, the system is less prone to sudden temporal spikes. 
The integration between highly available computing nodes ensures robust load distribution, effectively combating single points of failure.
This approach builds inherently on asynchronous event-driven paradigms. 


Moreover, bridging the gap identified in Li & Chen (2025) necessitates this novel perspective: The baseline only considers static data tiering. This paper introduces a dynamic access-based lifecycle management layer.

As noted in Mitchell (2025), scalability remains paramount. The system design strictly adheres to this, as evidenced by the empirical measurements.

The design isolates the core variables.


Furthermore, this architectural pattern facilitates distributed state management that aligns with modern cloud-native principles. 
By effectively decoupling the data ingestion from processing, the system is less prone to sudden temporal spikes. 
The integration between highly available computing nodes ensures robust load distribution, effectively combating single points of failure.
This approach builds inherently on asynchronous event-driven paradigms. 


The design isolates the core variables.


Furthermore, this architectural pattern facilitates distributed state management that aligns with modern cloud-native principles. 
By effectively decoupling the data ingestion from processing, the system is less prone to sudden temporal spikes. 
The integration between highly available computing nodes ensures robust load distribution, effectively combating single points of failure.
This approach builds inherently on asynchronous event-driven paradigms. 


The design isolates the core variables.

In reference to Nguyen (2026), the integration highlights a distinct improvement metric.


Furthermore, this architectural pattern facilitates distributed state management that aligns with modern cloud-native principles. 
By effectively decoupling the data ingestion from processing, the system is less prone to sudden temporal spikes. 
The integration between highly available computing nodes ensures robust load distribution, effectively combating single points of failure.
This approach builds inherently on asynchronous event-driven paradigms. 


The design isolates the core variables.


Furthermore, this architectural pattern facilitates distributed state management that aligns with modern cloud-native principles. 
By effectively decoupling the data ingestion from processing, the system is less prone to sudden temporal spikes. 
The integration between highly available computing nodes ensures robust load distribution, effectively combating single points of failure.
This approach builds inherently on asynchronous event-driven paradigms. 
